\section{The specified original fraction charts and their actual form bases}

\begin{lemma}[The specified generators cover their original Rees spectrum]
\label{lem:resume-specified-generator-basic-cover}
\uses{thm:rees-polynomial-surjective}
\lean{EnvelopeProperty.BlowupHomogeneousTautologicalLine.iSup_specifiedGeneratorChartOpen}
\leanok
Let $A$ be a commutative ring and let $g_0,\ldots,g_m$ generate an ideal $I$.
The actual basic opens $D_+(g_aT)$ cover $\operatorname{Proj}\mathcal R(I)$,
where $\mathcal R(I)\subset A[T]$ is the original Rees algebra.
\end{lemma}
\begin{proof}
The original graded polynomial map sending the $a$th variable to $g_aT$
is surjective because the given elements generate $I$. Hence these same
degree-one elements generate the whole Rees algebra over $A$.
They also generate it over its actual degree-zero subring: the subalgebra
over that subring contains every original coefficient through its degree-zero
Rees constant, and contains all the specified degree-one generators.
Induction on the original algebra generation therefore puts every Rees element
in this subalgebra. Basic opens of homogeneous positive-degree algebra generators
over degree zero cover the native projective spectrum. Indeed, a homogeneous
prime containing every such generator would contain the irrelevant ideal,
contrary to being a point of that projective spectrum.
\end{proof}

\begin{lemma}[The nonzero specified generators still cover]
\label{lem:resume-specified-nonzero-generator-basic-cover}
\uses{lem:resume-specified-generator-basic-cover}
\lean{EnvelopeProperty.BlowupHomogeneousTautologicalLine.iSup_specifiedNonzeroGeneratorChartOpen}
\leanok
In the preceding situation, the subfamily $D_+(g_aT)$ with $g_a\ne0$ still
covers the same actual Rees projective spectrum.
\end{lemma}
\begin{proof}
For a point of the projective spectrum, choose a specified generator whose
basic open contains it using the preceding lemma. If this generator were zero,
its degree-one Rees element would also be zero, and the corresponding basic open
would be empty. Thus the chosen generator is nonzero. This proves coverage by
the indicated subfamily. In particular, for the zero ideal all generators are
zero and this argument proves that the projective spectrum is empty.
\end{proof}

\begin{theorem}[The actual specified-generator fraction open cover]
\label{thm:resume-specified-generator-fraction-cover}
\uses{lem:resume-specified-nonzero-generator-basic-cover,lem:affine-blowup-fraction-chart-range}
\lean{EnvelopeProperty.BlowupHomogeneousTautologicalLine.specifiedGeneratorFractionOpenCover}
\leanok
Assume additionally that $A$ is a domain. The actual native morphisms
\[
 v_a:\operatorname{Spec}A[I/g_a]\longrightarrow
          \operatorname{Bl}_I(\operatorname{Spec}A),\qquad g_a\ne0,
\]
form a scheme open cover, with exactly these original finite-family indices
and actual fraction-chart morphisms.
\end{theorem}
\begin{proof}
Each displayed map is the native fraction-chart open immersion. Its open range
is exactly $D_+(g_aT)$, by the actual homogeneous-chart ring isomorphism and
native basic-open immersion. The preceding lemma shows that these ranges cover
the whole blowup. The schemes and maps therefore form the asserted scheme open
cover. If all the generators are zero, its index set and the native blowup are
both empty, and the same construction gives the empty open cover.
\end{proof}

\begin{lemma}[The whole original ideal on an actual fraction chart]
\label{lem:resume-original-fraction-chart-comap}
\uses{lem:blowup-exceptional-ideal-chart}
\lean{EnvelopeProperty.BlowupHomogeneousTautologicalLine.originalFractionChart_comap}
\leanok
Let $D$ be an ideal sheaf on a scheme $X$, let $U\subset X$ be affine with
coordinate domain $A$, and choose $0\ne g\in I=D(U)$.
Put $S=A[I/g]$, and let
\[
 f:\operatorname{Spec}S\xrightarrow{v}
      \operatorname{Bl}_I(\operatorname{Spec}A)
      \longrightarrow\operatorname{Spec}A\longrightarrow X
\]
be the composite of the actual native fraction chart, actual blowup projection,
and original affine spectrum map. Then the entire original ideal $D$ pulled
back by this composite is the ideal sheaf of the principal ideal $(\bar g)$
in $S$, where $\bar g$ is the image of the same original $g$.
\end{lemma}
\begin{proof}
The original affine spectrum comparison carries the whole ideal $D$ to the
ideal sheaf associated with $I=D(U)$. Ideal-sheaf pullback composes, so its
further pullback to the native blowup is the actual exceptional ideal, and its
pullback by $v$ is the actual exceptional ideal on that fraction chart.
The genuine chart structure triangle identifies this ideal with the extension
$IS$. Each $a\in I$ satisfies $\bar a=\bar g(a/g)$ in the actual subalgebra
$S\subset\operatorname{Frac}A$, so $IS\subset(\bar g)$. Conversely $g\in I$
gives $(\bar g)\subset IS$. This proves equality of the entire ideals and hence
of their associated ideal sheaves; no radical or support replacement occurs.
\end{proof}

\begin{theorem}[The original lifted form is a basis on its actual fraction chart]
\label{thm:resume-original-form-actual-fraction-basis}
\uses{lem:resume-original-fraction-chart-comap,lem:blowup-chart-equation-regular,thm:resume-original-lifted-form-regular-chart-basis}
\lean{EnvelopeProperty.BlowupHomogeneousTautologicalLine.liftedFormSection_isIso_on_originalFractionChart}
\leanok
Fix a commutative ring $R$, an original morphism $j:X\to\mathbf P_R^n$,
a homogeneous ideal $J$, and an original homogeneous $F\in J$ of degree
$d\in\mathbb N$. Let $U\subset j^{-1}D_+(X_i)$ be an affine open whose
coordinate ring $A$ is a domain. Write
\[
 D=j^{-1}\widetilde J,\qquad I=D(U),\qquad
 g=j^*(F/X_i^d)\in A,
\]
and assume $g\ne0$. Let $S=A[I/g]$, and let $f:\operatorname{Spec}S\to X$
be the actual composite from the preceding lemma. The actual original lift
\[
 \beta_{F,f}:\mathcal O_{\operatorname{Spec}S}\longrightarrow
          f^*j^*\mathcal O(d)\otimes D\mathcal O_{\operatorname{Spec}S}
\]
constructed from that same original $F$ and that same original $f$ is an
isomorphism. No regularity, principal-equation, section-matching, or
local-basis certificate is supplied.
\end{theorem}
\begin{proof}
Membership $F\in J$ implies $g\in I$: all original dehomogenized members of
$J$ generate the entire original coordinate ideal, and extension along the
original $j$ gives the whole ideal $D(U)$. Thus the specified native fraction
chart is defined from these original objects.

The fraction-chart structure triangle identifies $f$ with the original affine
spectrum map preceded by $\operatorname{Spec}(A\to S)$. Functoriality of the
actual section-ring morphisms therefore computes the original coefficient of
$F$ after pullback by $f$ as the image of the same $g$ under
\[
 A\longrightarrow S\longrightarrow\Gamma(\operatorname{Spec}S,\mathcal O).
\]
The final arrow is the canonical spectrum section-ring isomorphism, whose
naturality gives equality of entire coefficients. The original coordinate
containment follows by pulling back the given containment of $U$.

The preceding lemma identifies the entire original ideal $D\mathcal O_S$
with $(\bar g)$. The element $\bar g$ is nonzero in the actual subring
$S\subset\operatorname{Frac}A$, so it is a non-zero-divisor. Transport through
the canonical section-ring isomorphism supplies the actual principal regular
ideal equation on the whole affine spectrum. Apply the original regular-chart
basis theorem to these internally derived data and the same original lifted
form, using the whole affine spectrum as its affine open. It proves that the
original lift becomes an isomorphism after the canonical affine spectrum
comparison of this affine scheme. That comparison is itself an isomorphism,
so its module pullback reflects isomorphisms. The original lift is therefore
an isomorphism on the actual fraction chart itself.
\end{proof}

\paragraph{The remaining whole-line comparison.}
The intended consumer is the comparison between the original projective
blowup map's pulled $\mathcal O(1)$ and the original
$\mu^*\mathcal O(d)\otimes D\mathcal O_{\operatorname{Bl}_D X}$.
Its remaining construction must retain the canonical open-comap comparison of
the original lifts, the original projective coordinate ratios on the specified
fraction cover, and every original section when gluing. Passing to the actual
normalized simultaneous blowup, constructing a fixed coupled ample line, and
retaining its $\mathbf G_m$ action are further consumers. The preceding fraction
chart results are the stated local producers for this route.
